% LaTeX rendering of the RISEx Word template's stated page specification.
% This is an author-supplied adaptation, not an official conference template.
\documentclass[11pt,letterpaper]{article}
\usepackage[top=20mm,left=15mm,right=15mm,bottom=15mm,headheight=10pt,headsep=6mm]{geometry}
\usepackage{fontspec}
\IfFontExistsTF{Times New Roman}{
  \setmainfont{Times New Roman}
}{
  \setmainfont{TeX Gyre Termes}
  \PackageWarningNoLine{risex}{Times New Roman unavailable. Termes is a proofing fallback; check the official font requirement before submission}
}
\usepackage{amsmath,amssymb}
\usepackage{unicode-math}
\setmathfont{texgyretermes-math.otf}
\usepackage{microtype}
\usepackage{multicol}
\usepackage{array,booktabs,tabularx}
\usepackage{fancyhdr}
\usepackage[hidelinks]{hyperref}
\hypersetup{pdftitle={Neurologically Motivated Self-Supervised Learning
for Bilateral Gait
Geometry},pdfauthor={Alexander Mui and Penelope Inouye},pdfsubject={RISEx 2026 one-page research draft}}
\pagestyle{fancy}
\fancyhf{}
\fancyhead[C]{\fontsize{8}{9}\selectfont Robotics \& Intelligent Systems Expo (RISEx 2026), Nov. 12-13, 2026, Edmonton, Canada}
\renewcommand{\headrulewidth}{0pt}
\setlength{\columnsep}{7.5mm}
\setlength{\parindent}{0pt}
\setlength{\parskip}{3pt}
\setlength{\multicolsep}{6pt}
\setlength{\abovedisplayskip}{5pt}
\setlength{\belowdisplayskip}{5pt}
\setlength{\abovedisplayshortskip}{3pt}
\setlength{\belowdisplayshortskip}{3pt}
\setlength{\emergencystretch}{1em}
\newcommand{\risexheading}[1]{\par\vspace{4pt}{\bfseries\underline{#1}}\par\nobreak\vspace{1pt}}
\providecommand{\tightlist}{\setlength{\itemsep}{0pt}\setlength{\parskip}{0pt}}
\begin{document}
\fontsize{11}{12.5}\selectfont
{\centering
  {\bfseries\MakeUppercase{Neurologically Motivated Self-Supervised}\\
   \MakeUppercase{Learning for Bilateral Gait Geometry}\par}
  Alexander Mui\textsuperscript{*}, Penelope Inouye\par
  Computer Science \& Engineering, ASDRP, Fremont, CA, USA\par
  \textsuperscript{*}Corresponding author: Alexander Mui; email: \href{mailto:alexander.mui@students.asdrp.org}{alexander.mui@students.asdrp.org}\par
}
\begin{multicols}{2}
\risexheading{INTRODUCTION}

Stroke-related walking asymmetry and asymmetric arm swing in Parkinson's
disease {[}1,2{]} motivate side-sensitive representations for
human-movement understanding and rehabilitation robotics. We test
whether a joint-embedding predictive architecture (JEPA), which predicts
hidden features rather than coordinates, improves recovery of a signed
bilateral speed contrast. Our exploratory null hypothesis is no
improvement over matched initialization.

\risexheading{MATERIALS AND METHODS}

We use 625 quality-checked GAVD pose clips from 93 source videos
{[}3{]}. For five joint pairs (shoulders, knees, ankles, heels and foot
tips), the target is

\[
y=\frac{1}{5}\sum_{k=1}^{5}\frac{m_{L,k}-m_{R,k}}{m_{L,k}+m_{R,k}+10^{-8}}.
\]

Here \(m\) is median speed in body-normalized pose coordinates, using
original timestamps and at least eight jointly observed transitions per
pair. Reflection swaps sides and gives \(y(Mx)=-y(x)\). This
pose-derived measure has no independent clinical validation.

Clip-local centering, scaling, short-gap filling and resampling produce
\(64\times33\times3\) inputs. Four-frame patches retain a \(16\times33\)
time--joint grid; masking zeroes features without removing positions.
Joint identity is preserved, while resampling changes the original
clock.

Our S-JEPA-inspired model {[}4{]} uses 96-channel features, a four-layer
encoder and two-layer predictor. Masked cross-entropy matches
teacher-channel distributions; the teacher receives no gradients and
follows exponential averaging (0.999). A 0.05-weighted VICReg term
{[}5{]} encourages agreement and nonconstant features pooled over 12
gait joints across two complete views. Neither target scores nor
neurological labels enter pretraining. Two motion-based masks and a
connected-region mask have count-matched random controls within each
mask family: five conditions, five folds and five seeds give 125 runs,
each with 1,200 updates and batches of 20. Paired conditions share
initialization, source draws and views.

Five outer folds exclude entire source videos from pretraining and
readout fitting. All encoder comparisons use the same 2,890-feature
summary of bilateral means, temporal variation and observation support.
A ridge readout (linear regression with shrinkage) uses three inner
source-separated folds for training-only scaling and penalty selection.
Held-out predictions are pooled across outer folds; sources receive
equal total weight, and scores are averaged across seeds. Paired 95\%
intervals use 2,000 whole-source bootstrap draws with fitted models
fixed.

\risexheading{RESULTS AND DISCUSSION}

Training reduced bilateral readout accuracy across all five conditions
(Table 1). For the motion-family random control, trained-minus-initial
\(\Delta R^2=-0.109\) {[}−0.170, −0.041{]}. All five marginal intervals
lie below zero; structured-mask comparisons with their random controls
span zero.

\begin{minipage}{\linewidth}
{\fontsize{9}{10.5}\selectfont \emph{Table 1. Source-balanced held-out scores. Higher \(R^2\) and lower
mean absolute error (MAE) are better; ranges cover five trained
conditions.}\par}
\vspace{3pt}

\setlength{\tabcolsep}{2pt}
\begin{tabularx}{\linewidth}{@{}>{\raggedright\arraybackslash}Xrr@{}}
\toprule
Predictor & \(R^2\) & MAE \\
\midrule
Training-source mean & −0.011 & 0.0462 \\
Matched initial encoder & 0.223 & 0.0415 \\
Trained teachers & 0.101--0.114 & 0.0436--0.0440 \\
\bottomrule
\end{tabularx}
\end{minipage}

Correct-clip teacher targets yield lower feature error than targets from
another source in 375/375 trained checks, versus 33/75 at
initialization. These repeated fold/seed/mask checks establish clip
correspondence, which can also reflect viewpoint or missingness. They do
not isolate movement prediction.

Original and preprocessed target calculations agree in sign on 70.4\% of
source-weighted comparisons (623 clips, 92 sources); 49/125 teacher
readouts select the largest tested ridge penalty. These observations
limit interpretation of the readout deficit. Analyses are exploratory;
intervals omit retraining and selection uncertainty, and video
separation does not verify participant independence.

\risexheading{CONCLUSIONS}

Under this recipe, clip-specific hidden-feature matching coexists with
poorer recovery of bilateral movement by the ridge readout. For
human-movement models, the result supports testing a fixed geometric
target alongside the training objective. Time-preserving inputs and a
wider readout-penalty search should be evaluated before drawing
conclusions about suitability for rehabilitation or future-motion
prediction.

\risexheading{REFERENCES}

{[}1{]} Patterson KK et al.~\emph{Arch Phys Med Rehabil} 89:304--310,
2008. {[}2{]} Lewek MD et al.~\emph{Gait Posture} 31:256--260, 2010.
{[}3{]} Ranjan R et al.~\emph{IEEE Access} 13:45321--45339, 2025.
{[}4{]} Abdelfattah M, Alahi A. \emph{ECCV 2024}, LNCS 15090:367--384.
{[}5{]} Bardes A et al.~\emph{ICLR}, 2022; arXiv:2105.04906.
\end{multicols}
\end{document}
